\begin{example}
\label{example-oka-family-bound-cardinality}
Let $R$ be a ring.
Let $\kappa$ be an infinite cardinal.
The family of ideals which can be generated by at most $\kappa$ elements
is an Oka family. To verify this, suppose $(I:a)$ has a
generating family $(u_\lambda)_{\lambda\in L}$ and $(I,a)$ has
a generating family $(v_\mu)_{\mu\in P}$, with $|L|,|P|\leq\kappa$.
Write $v_\mu=b_\mu+c_\mu a$, with $b_\mu\in I$ and $c_\mu\in R$.
Then $(I,a)=(a,b_\mu\mid\mu\in P)$.
For $x\in I$, an expression
$x=ra+\sum_{\mu\in P_0}s_\mu b_\mu$ with $P_0$ finite gives
$ra\in I$, hence $r\in(I:a)$. Writing
$r=\sum_{\lambda\in L_0}t_\lambda u_\lambda$ with $L_0$ finite yields
$$
x=\sum_{\lambda\in L_0}t_\lambda a u_\lambda+
\sum_{\mu\in P_0}s_\mu b_\mu.
$$
Thus $I=(a u_\lambda,b_\mu\mid\lambda\in L,\mu\in P)$, with
at most $\kappa+\kappa=\kappa$ generators. The ideal $R$ is
generated by $1$, so it belongs to this family as well.
\end{example}